\section{余下账本事件导航：eastern-han}

\label{sec:remaining-event-anchors-eastern-han}

本节为尚未在正文设置目标的账本记录建立直接跳转。锚点显示可核日期与事件名称；其解释仍应回到账本所列来源、置信提示及既有正文的区域和材料边界。

\subsection{事件导航与来源边界}

\eventanchor{EH-LEDGER-007-6}{建武十一年（35年）·“光武帝下诏限制将战乱中掠卖者继续为奴婢，并令核实释放”的异源材料与影响范围的复核}

\eventanchor{EH-LEDGER-008-1}{建武十二年（36年）·“吴汉等攻入成都，公孙述战死，成家政权灭亡”的前期动员与资源准备}

\eventanchor{EH-LEDGER-008-2}{建武十二年（36年）·“吴汉等攻入成都，公孙述战死，成家政权灭亡”的命令、联盟或地方响应}

\eventanchor{EH-LEDGER-008-3}{建武十二年（36年）·吴汉等攻入成都，公孙述战死，成家政权灭亡}

\eventanchor{EH-LEDGER-008-4}{建武十二年（36年）·“吴汉等攻入成都，公孙述战死，成家政权灭亡”的执行中的空间与人员处置}

\eventanchor{EH-LEDGER-008-5}{建武十二年（36年）·“吴汉等攻入成都，公孙述战死，成家政权灭亡”的后续制度或军政调整}

\eventanchor{EH-LEDGER-008-6}{建武十二年（36年）·“吴汉等攻入成都，公孙述战死，成家政权灭亡”的异源材料与影响范围的复核}

\eventanchor{EH-LEDGER-009-1}{建武十五年（39年）·“光武帝命郡国清查垦田、户口，因地方豪强隐匿与吏治争议引发纠纷”的前期动员与资源准备}

\eventanchor{EH-LEDGER-009-2}{建武十五年（39年）·“光武帝命郡国清查垦田、户口，因地方豪强隐匿与吏治争议引发纠纷”的命令、联盟或地方响应}

\eventanchor{EH-LEDGER-009-3}{建武十五年（39年）·光武帝命郡国清查垦田、户口，因地方豪强隐匿与吏治争议引发纠纷}

\eventanchor{EH-LEDGER-009-4}{建武十五年（39年）·“光武帝命郡国清查垦田、户口，因地方豪强隐匿与吏治争议引发纠纷”的执行中的空间与人员处置}

\eventanchor{EH-LEDGER-009-5}{建武十五年（39年）·“光武帝命郡国清查垦田、户口，因地方豪强隐匿与吏治争议引发纠纷”的后续制度或军政调整}

\eventanchor{EH-LEDGER-009-6}{建武十五年（39年）·“光武帝命郡国清查垦田、户口，因地方豪强隐匿与吏治争议引发纠纷”的异源材料与影响范围的复核}

\eventanchor{EH-LEDGER-010-1}{建武十六年（40年）·“徵侧、徵贰在交趾起兵，攻取多处郡县，东汉岭南统治受挑战”的前期动员与资源准备}

\eventanchor{EH-LEDGER-010-2}{建武十六年（40年）·“徵侧、徵贰在交趾起兵，攻取多处郡县，东汉岭南统治受挑战”的命令、联盟或地方响应}

\eventanchor{EH-LEDGER-010-3}{建武十六年（40年）·徵侧、徵贰在交趾起兵，攻取多处郡县，东汉岭南统治受挑战}

\eventanchor{EH-LEDGER-010-4}{建武十六年（40年）·“徵侧、徵贰在交趾起兵，攻取多处郡县，东汉岭南统治受挑战”的执行中的空间与人员处置}

\eventanchor{EH-LEDGER-010-5}{建武十六年（40年）·“徵侧、徵贰在交趾起兵，攻取多处郡县，东汉岭南统治受挑战”的后续制度或军政调整}

\eventanchor{EH-LEDGER-010-6}{建武十六年（40年）·“徵侧、徵贰在交趾起兵，攻取多处郡县，东汉岭南统治受挑战”的异源材料与影响范围的复核}

\eventanchor{EH-LEDGER-011-1}{建武十九年（43年）·“马援平定交趾，徵氏集团覆灭，东汉重建岭南郡县控制”的前期动员与资源准备}

\eventanchor{EH-LEDGER-011-2}{建武十九年（43年）·“马援平定交趾，徵氏集团覆灭，东汉重建岭南郡县控制”的命令、联盟或地方响应}

\eventanchor{EH-LEDGER-011-3}{建武十九年（43年）·马援平定交趾，徵氏集团覆灭，东汉重建岭南郡县控制}

\eventanchor{EH-LEDGER-011-4}{建武十九年（43年）·“马援平定交趾，徵氏集团覆灭，东汉重建岭南郡县控制”的执行中的空间与人员处置}

\eventanchor{EH-LEDGER-011-5}{建武十九年（43年）·“马援平定交趾，徵氏集团覆灭，东汉重建岭南郡县控制”的后续制度或军政调整}

\eventanchor{EH-LEDGER-011-6}{建武十九年（43年）·“马援平定交趾，徵氏集团覆灭，东汉重建岭南郡县控制”的异源材料与影响范围的复核}

\eventanchor{EH-LEDGER-012-1}{建武二十四年（48年）·“南匈奴日逐王比归附汉廷，南北匈奴分裂，汉在边郡安置其部众”的前期动员与资源准备}

\eventanchor{EH-LEDGER-012-2}{建武二十四年（48年）·“南匈奴日逐王比归附汉廷，南北匈奴分裂，汉在边郡安置其部众”的命令、联盟或地方响应}

\eventanchor{EH-LEDGER-012-3}{建武二十四年（48年）·南匈奴日逐王比归附汉廷，南北匈奴分裂，汉在边郡安置其部众}

\eventanchor{EH-LEDGER-012-4}{建武二十四年（48年）·“南匈奴日逐王比归附汉廷，南北匈奴分裂，汉在边郡安置其部众”的执行中的空间与人员处置}

\eventanchor{EH-LEDGER-012-5}{建武二十四年（48年）·“南匈奴日逐王比归附汉廷，南北匈奴分裂，汉在边郡安置其部众”的后续制度或军政调整}

\eventanchor{EH-LEDGER-012-6}{建武二十四年（48年）·“南匈奴日逐王比归附汉廷，南北匈奴分裂，汉在边郡安置其部众”的异源材料与影响范围的复核}

\eventanchor{EH-LEDGER-013-1}{建武中元二年（57年）·“光武帝去世，太子刘庄即位，是为明帝”的前期动员与资源准备}

\eventanchor{EH-LEDGER-013-2}{建武中元二年（57年）·“光武帝去世，太子刘庄即位，是为明帝”的命令、联盟或地方响应}

\eventanchor{EH-LEDGER-013-3}{建武中元二年（57年）·光武帝去世，太子刘庄即位，是为明帝}

\eventanchor{EH-LEDGER-013-4}{建武中元二年（57年）·“光武帝去世，太子刘庄即位，是为明帝”的执行中的空间与人员处置}

\eventanchor{EH-LEDGER-013-5}{建武中元二年（57年）·“光武帝去世，太子刘庄即位，是为明帝”的后续制度或军政调整}

\eventanchor{EH-LEDGER-013-6}{建武中元二年（57年）·“光武帝去世，太子刘庄即位，是为明帝”的异源材料与影响范围的复核}

\eventanchor{EH-LEDGER-014-1}{永平八年（65年）·“楚王刘英因被告谋反获减死并被徙，其诏书涉及黄老与佛教供养的早期记录”的前期动员与资源准备}

\eventanchor{EH-LEDGER-014-2}{永平八年（65年）·“楚王刘英因被告谋反获减死并被徙，其诏书涉及黄老与佛教供养的早期记录”的命令、联盟或地方响应}

\eventanchor{EH-LEDGER-014-3}{永平八年（65年）·楚王刘英因被告谋反获减死并被徙，其诏书涉及黄老与佛教供养的早期记录}

\eventanchor{EH-LEDGER-014-4}{永平八年（65年）·“楚王刘英因被告谋反获减死并被徙，其诏书涉及黄老与佛教供养的早期记录”的执行中的空间与人员处置}

\eventanchor{EH-LEDGER-014-5}{永平八年（65年）·“楚王刘英因被告谋反获减死并被徙，其诏书涉及黄老与佛教供养的早期记录”的后续制度或军政调整}

\eventanchor{EH-LEDGER-014-6}{永平八年（65年）·“楚王刘英因被告谋反获减死并被徙，其诏书涉及黄老与佛教供养的早期记录”的异源材料与影响范围的复核}

\eventanchor{EH-LEDGER-015-1}{永平十六年（73年）·“窦固等出击北匈奴，取伊吾卢；班超奉使鄯善，东汉重启西域经营”的前期动员与资源准备}

\eventanchor{EH-LEDGER-015-2}{永平十六年（73年）·“窦固等出击北匈奴，取伊吾卢；班超奉使鄯善，东汉重启西域经营”的命令、联盟或地方响应}

\eventanchor{EH-LEDGER-015-3}{永平十六年（73年）·窦固等出击北匈奴，取伊吾卢；班超奉使鄯善，东汉重启西域经营}

\eventanchor{EH-LEDGER-015-4}{永平十六年（73年）·“窦固等出击北匈奴，取伊吾卢；班超奉使鄯善，东汉重启西域经营”的执行中的空间与人员处置}

\eventanchor{EH-LEDGER-015-5}{永平十六年（73年）·“窦固等出击北匈奴，取伊吾卢；班超奉使鄯善，东汉重启西域经营”的后续制度或军政调整}

\eventanchor{EH-LEDGER-015-6}{永平十六年（73年）·“窦固等出击北匈奴，取伊吾卢；班超奉使鄯善，东汉重启西域经营”的异源材料与影响范围的复核}

\eventanchor{EH-LEDGER-016-1}{永平十八年（75年）·“北匈奴攻车师，西域都护陈睦遇害，东汉一度撤出西域据点”的前期动员与资源准备}

\eventanchor{EH-LEDGER-016-2}{永平十八年（75年）·“北匈奴攻车师，西域都护陈睦遇害，东汉一度撤出西域据点”的命令、联盟或地方响应}

\eventanchor{EH-LEDGER-016-3}{永平十八年（75年）·北匈奴攻车师，西域都护陈睦遇害，东汉一度撤出西域据点}

\eventanchor{EH-LEDGER-016-4}{永平十八年（75年）·“北匈奴攻车师，西域都护陈睦遇害，东汉一度撤出西域据点”的执行中的空间与人员处置}

\eventanchor{EH-LEDGER-016-5}{永平十八年（75年）·“北匈奴攻车师，西域都护陈睦遇害，东汉一度撤出西域据点”的后续制度或军政调整}

\eventanchor{EH-LEDGER-016-6}{永平十八年（75年）·“北匈奴攻车师，西域都护陈睦遇害，东汉一度撤出西域据点”的异源材料与影响范围的复核}

\eventanchor{EH-LEDGER-017-1}{永平十八年（75年）·“明帝去世，太子刘炟即位，是为章帝”的前期动员与资源准备}

\eventanchor{EH-LEDGER-017-2}{永平十八年（75年）·“明帝去世，太子刘炟即位，是为章帝”的命令、联盟或地方响应}

\eventanchor{EH-LEDGER-017-3}{永平十八年（75年）·明帝去世，太子刘炟即位，是为章帝}

\eventanchor{EH-LEDGER-017-4}{永平十八年（75年）·“明帝去世，太子刘炟即位，是为章帝”的执行中的空间与人员处置}

\eventanchor{EH-LEDGER-017-5}{永平十八年（75年）·“明帝去世，太子刘炟即位，是为章帝”的后续制度或军政调整}

\subsection{事件导航与来源边界}

\eventanchor{EH-LEDGER-017-6}{永平十八年（75年）·“明帝去世，太子刘炟即位，是为章帝”的异源材料与影响范围的复核}

\eventanchor{EH-LEDGER-018-1}{建初元年（76年）·“西羌在陇西、金城等地反叛，东汉长期羌乱由此加剧”的前期动员与资源准备}

\eventanchor{EH-LEDGER-018-2}{建初元年（76年）·“西羌在陇西、金城等地反叛，东汉长期羌乱由此加剧”的命令、联盟或地方响应}

\eventanchor{EH-LEDGER-018-3}{建初元年（76年）·西羌在陇西、金城等地反叛，东汉长期羌乱由此加剧}

\eventanchor{EH-LEDGER-018-4}{建初元年（76年）·“西羌在陇西、金城等地反叛，东汉长期羌乱由此加剧”的执行中的空间与人员处置}

\eventanchor{EH-LEDGER-018-5}{建初元年（76年）·“西羌在陇西、金城等地反叛，东汉长期羌乱由此加剧”的后续制度或军政调整}

\eventanchor{EH-LEDGER-018-6}{建初元年（76年）·“西羌在陇西、金城等地反叛，东汉长期羌乱由此加剧”的异源材料与影响范围的复核}

\eventanchor{EH-LEDGER-019-1}{建初四年（79年）·“章帝召集儒者在白虎观讨论五经异同，形成《白虎通》所保存的经义框架”的前期动员与资源准备}

\eventanchor{EH-LEDGER-019-2}{建初四年（79年）·“章帝召集儒者在白虎观讨论五经异同，形成《白虎通》所保存的经义框架”的命令、联盟或地方响应}

\eventanchor{EH-LEDGER-019-3}{建初四年（79年）·章帝召集儒者在白虎观讨论五经异同，形成《白虎通》所保存的经义框架}

\eventanchor{EH-LEDGER-019-4}{建初四年（79年）·“章帝召集儒者在白虎观讨论五经异同，形成《白虎通》所保存的经义框架”的执行中的空间与人员处置}

\eventanchor{EH-LEDGER-019-5}{建初四年（79年）·“章帝召集儒者在白虎观讨论五经异同，形成《白虎通》所保存的经义框架”的后续制度或军政调整}

\eventanchor{EH-LEDGER-019-6}{建初四年（79年）·“章帝召集儒者在白虎观讨论五经异同，形成《白虎通》所保存的经义框架”的异源材料与影响范围的复核}

\eventanchor{EH-LEDGER-020-1}{章和二年（88年）·“章帝去世，皇子刘肇即位，是为和帝，窦太后临朝”的前期动员与资源准备}

\eventanchor{EH-LEDGER-020-2}{章和二年（88年）·“章帝去世，皇子刘肇即位，是为和帝，窦太后临朝”的命令、联盟或地方响应}

\eventanchor{EH-LEDGER-020-3}{章和二年（88年）·章帝去世，皇子刘肇即位，是为和帝，窦太后临朝}

\eventanchor{EH-LEDGER-020-4}{章和二年（88年）·“章帝去世，皇子刘肇即位，是为和帝，窦太后临朝”的执行中的空间与人员处置}

\eventanchor{EH-LEDGER-020-5}{章和二年（88年）·“章帝去世，皇子刘肇即位，是为和帝，窦太后临朝”的后续制度或军政调整}

\eventanchor{EH-LEDGER-020-6}{章和二年（88年）·“章帝去世，皇子刘肇即位，是为和帝，窦太后临朝”的异源材料与影响范围的复核}

\eventanchor{EH-LEDGER-021-1}{永元元年（89年）·“窦宪等北出稽落山，大破北匈奴，立燕然山铭纪功”的前期动员与资源准备}

\eventanchor{EH-LEDGER-021-2}{永元元年（89年）·“窦宪等北出稽落山，大破北匈奴，立燕然山铭纪功”的命令、联盟或地方响应}

\eventanchor{EH-LEDGER-021-3}{永元元年（89年）·窦宪等北出稽落山，大破北匈奴，立燕然山铭纪功}

\eventanchor{EH-LEDGER-021-4}{永元元年（89年）·“窦宪等北出稽落山，大破北匈奴，立燕然山铭纪功”的执行中的空间与人员处置}

\eventanchor{EH-LEDGER-021-5}{永元元年（89年）·“窦宪等北出稽落山，大破北匈奴，立燕然山铭纪功”的后续制度或军政调整}

\eventanchor{EH-LEDGER-021-6}{永元元年（89年）·“窦宪等北出稽落山，大破北匈奴，立燕然山铭纪功”的异源材料与影响范围的复核}

\eventanchor{EH-LEDGER-022-1}{永元三年（91年）·“班超受命为西域都护，东汉在多处绿洲重新建立使节与军事网络”的前期动员与资源准备}

\eventanchor{EH-LEDGER-022-2}{永元三年（91年）·“班超受命为西域都护，东汉在多处绿洲重新建立使节与军事网络”的命令、联盟或地方响应}

\eventanchor{EH-LEDGER-022-3}{永元三年（91年）·班超受命为西域都护，东汉在多处绿洲重新建立使节与军事网络}

\eventanchor{EH-LEDGER-022-4}{永元三年（91年）·“班超受命为西域都护，东汉在多处绿洲重新建立使节与军事网络”的执行中的空间与人员处置}

\eventanchor{EH-LEDGER-022-5}{永元三年（91年）·“班超受命为西域都护，东汉在多处绿洲重新建立使节与军事网络”的后续制度或军政调整}

\eventanchor{EH-LEDGER-022-6}{永元三年（91年）·“班超受命为西域都护，东汉在多处绿洲重新建立使节与军事网络”的异源材料与影响范围的复核}

\eventanchor{EH-LEDGER-023-1}{永元四年（92年）·“和帝联同宦官郑众诛除窦宪集团，结束窦太后家族的军政主导”的前期动员与资源准备}

\eventanchor{EH-LEDGER-023-2}{永元四年（92年）·“和帝联同宦官郑众诛除窦宪集团，结束窦太后家族的军政主导”的命令、联盟或地方响应}

\eventanchor{EH-LEDGER-023-3}{永元四年（92年）·和帝联同宦官郑众诛除窦宪集团，结束窦太后家族的军政主导}

\eventanchor{EH-LEDGER-023-4}{永元四年（92年）·“和帝联同宦官郑众诛除窦宪集团，结束窦太后家族的军政主导”的执行中的空间与人员处置}

\eventanchor{EH-LEDGER-023-5}{永元四年（92年）·“和帝联同宦官郑众诛除窦宪集团，结束窦太后家族的军政主导”的后续制度或军政调整}

\eventanchor{EH-LEDGER-023-6}{永元四年（92年）·“和帝联同宦官郑众诛除窦宪集团，结束窦太后家族的军政主导”的异源材料与影响范围的复核}

\eventanchor{EH-LEDGER-024-1}{永元九年（97年）·“班超遣甘英西使，至安息属地而未能到达大秦”的前期动员与资源准备}

\eventanchor{EH-LEDGER-024-2}{永元九年（97年）·“班超遣甘英西使，至安息属地而未能到达大秦”的命令、联盟或地方响应}

\eventanchor{EH-LEDGER-024-3}{永元九年（97年）·班超遣甘英西使，至安息属地而未能到达大秦}

\eventanchor{EH-LEDGER-024-4}{永元九年（97年）·“班超遣甘英西使，至安息属地而未能到达大秦”的执行中的空间与人员处置}

\eventanchor{EH-LEDGER-024-5}{永元九年（97年）·“班超遣甘英西使，至安息属地而未能到达大秦”的后续制度或军政调整}

\eventanchor{EH-LEDGER-024-6}{永元九年（97年）·“班超遣甘英西使，至安息属地而未能到达大秦”的异源材料与影响范围的复核}

\eventanchor{EH-LEDGER-025-1}{永元十四年（102年）·“班超返洛阳后去世，东汉西域经营失去长期协调者”的前期动员与资源准备}

\eventanchor{EH-LEDGER-025-2}{永元十四年（102年）·“班超返洛阳后去世，东汉西域经营失去长期协调者”的命令、联盟或地方响应}

\eventanchor{EH-LEDGER-025-3}{永元十四年（102年）·班超返洛阳后去世，东汉西域经营失去长期协调者}

\eventanchor{EH-LEDGER-025-4}{永元十四年（102年）·“班超返洛阳后去世，东汉西域经营失去长期协调者”的执行中的空间与人员处置}

\eventanchor{EH-LEDGER-025-5}{永元十四年（102年）·“班超返洛阳后去世，东汉西域经营失去长期协调者”的后续制度或军政调整}

\eventanchor{EH-LEDGER-025-6}{永元十四年（102年）·“班超返洛阳后去世，东汉西域经营失去长期协调者”的异源材料与影响范围的复核}

\eventanchor{EH-LEDGER-026-1}{元兴元年（105年）·“和帝去世，襁褓中的刘隆即位，是为殇帝，邓太后临朝”的前期动员与资源准备}

\eventanchor{EH-LEDGER-026-2}{元兴元年（105年）·“和帝去世，襁褓中的刘隆即位，是为殇帝，邓太后临朝”的命令、联盟或地方响应}

\eventanchor{EH-LEDGER-026-3}{元兴元年（105年）·和帝去世，襁褓中的刘隆即位，是为殇帝，邓太后临朝}

\eventanchor{EH-LEDGER-026-4}{元兴元年（105年）·“和帝去世，襁褓中的刘隆即位，是为殇帝，邓太后临朝”的执行中的空间与人员处置}

\eventanchor{EH-LEDGER-026-5}{元兴元年（105年）·“和帝去世，襁褓中的刘隆即位，是为殇帝，邓太后临朝”的后续制度或军政调整}

\eventanchor{EH-LEDGER-026-6}{元兴元年（105年）·“和帝去世，襁褓中的刘隆即位，是为殇帝，邓太后临朝”的异源材料与影响范围的复核}

\eventanchor{EH-LEDGER-027-1}{延平元年（106年）·“殇帝去世，清河王刘祜即位，是为安帝，邓太后继续临朝”的前期动员与资源准备}

\eventanchor{EH-LEDGER-027-2}{延平元年（106年）·“殇帝去世，清河王刘祜即位，是为安帝，邓太后继续临朝”的命令、联盟或地方响应}

\eventanchor{EH-LEDGER-027-3}{延平元年（106年）·殇帝去世，清河王刘祜即位，是为安帝，邓太后继续临朝}

\eventanchor{EH-LEDGER-027-4}{延平元年（106年）·“殇帝去世，清河王刘祜即位，是为安帝，邓太后继续临朝”的执行中的空间与人员处置}

\eventanchor{EH-LEDGER-027-5}{延平元年（106年）·“殇帝去世，清河王刘祜即位，是为安帝，邓太后继续临朝”的后续制度或军政调整}

\subsection{事件导航与来源边界}

\eventanchor{EH-LEDGER-027-6}{延平元年（106年）·“殇帝去世，清河王刘祜即位，是为安帝，邓太后继续临朝”的异源材料与影响范围的复核}

\eventanchor{EH-LEDGER-028-1}{永初元年（107年）·“西羌在凉州、并州一带再起，东汉长期投入边郡军费与移民安置”的前期动员与资源准备}

\eventanchor{EH-LEDGER-028-2}{永初元年（107年）·“西羌在凉州、并州一带再起，东汉长期投入边郡军费与移民安置”的命令、联盟或地方响应}

\eventanchor{EH-LEDGER-028-3}{永初元年（107年）·西羌在凉州、并州一带再起，东汉长期投入边郡军费与移民安置}

\eventanchor{EH-LEDGER-028-4}{永初元年（107年）·“西羌在凉州、并州一带再起，东汉长期投入边郡军费与移民安置”的执行中的空间与人员处置}

\eventanchor{EH-LEDGER-028-5}{永初元年（107年）·“西羌在凉州、并州一带再起，东汉长期投入边郡军费与移民安置”的后续制度或军政调整}

\eventanchor{EH-LEDGER-028-6}{永初元年（107年）·“西羌在凉州、并州一带再起，东汉长期投入边郡军费与移民安置”的异源材料与影响范围的复核}

\eventanchor{EH-LEDGER-029-1}{元初三年（116年）·“邓太后去世，安帝开始亲政并贬逐部分邓氏成员”的前期动员与资源准备}

\eventanchor{EH-LEDGER-029-2}{元初三年（116年）·“邓太后去世，安帝开始亲政并贬逐部分邓氏成员”的命令、联盟或地方响应}

\eventanchor{EH-LEDGER-029-3}{元初三年（116年）·邓太后去世，安帝开始亲政并贬逐部分邓氏成员}

\eventanchor{EH-LEDGER-029-4}{元初三年（116年）·“邓太后去世，安帝开始亲政并贬逐部分邓氏成员”的执行中的空间与人员处置}

\eventanchor{EH-LEDGER-029-5}{元初三年（116年）·“邓太后去世，安帝开始亲政并贬逐部分邓氏成员”的后续制度或军政调整}

\eventanchor{EH-LEDGER-029-6}{元初三年（116年）·“邓太后去世，安帝开始亲政并贬逐部分邓氏成员”的异源材料与影响范围的复核}

\eventanchor{EH-LEDGER-030-1}{延光四年（125年）·“安帝去世，阎后先立北乡侯，宦官孙程等旋迎立济阴王刘保为顺帝”的前期动员与资源准备}

\eventanchor{EH-LEDGER-030-2}{延光四年（125年）·“安帝去世，阎后先立北乡侯，宦官孙程等旋迎立济阴王刘保为顺帝”的命令、联盟或地方响应}

\eventanchor{EH-LEDGER-030-3}{延光四年（125年）·安帝去世，阎后先立北乡侯，宦官孙程等旋迎立济阴王刘保为顺帝}

\eventanchor{EH-LEDGER-030-4}{延光四年（125年）·“安帝去世，阎后先立北乡侯，宦官孙程等旋迎立济阴王刘保为顺帝”的执行中的空间与人员处置}

\eventanchor{EH-LEDGER-030-5}{延光四年（125年）·“安帝去世，阎后先立北乡侯，宦官孙程等旋迎立济阴王刘保为顺帝”的后续制度或军政调整}

\eventanchor{EH-LEDGER-030-6}{延光四年（125年）·“安帝去世，阎后先立北乡侯，宦官孙程等旋迎立济阴王刘保为顺帝”的异源材料与影响范围的复核}

\eventanchor{EH-LEDGER-031-1}{建康元年（144年）·“顺帝去世，幼帝刘炳即位，是为冲帝，梁太后与梁冀临朝”的前期动员与资源准备}

\eventanchor{EH-LEDGER-031-2}{建康元年（144年）·“顺帝去世，幼帝刘炳即位，是为冲帝，梁太后与梁冀临朝”的命令、联盟或地方响应}

\eventanchor{EH-LEDGER-031-3}{建康元年（144年）·顺帝去世，幼帝刘炳即位，是为冲帝，梁太后与梁冀临朝}

\eventanchor{EH-LEDGER-031-4}{建康元年（144年）·“顺帝去世，幼帝刘炳即位，是为冲帝，梁太后与梁冀临朝”的执行中的空间与人员处置}

\eventanchor{EH-LEDGER-031-5}{建康元年（144年）·“顺帝去世，幼帝刘炳即位，是为冲帝，梁太后与梁冀临朝”的后续制度或军政调整}

\eventanchor{EH-LEDGER-031-6}{建康元年（144年）·“顺帝去世，幼帝刘炳即位，是为冲帝，梁太后与梁冀临朝”的异源材料与影响范围的复核}

\eventanchor{EH-LEDGER-032-1}{永嘉元年（145年）·“冲帝夭亡，梁太后立八岁的刘缵为质帝，梁冀继续掌政”的前期动员与资源准备}

\eventanchor{EH-LEDGER-032-2}{永嘉元年（145年）·“冲帝夭亡，梁太后立八岁的刘缵为质帝，梁冀继续掌政”的命令、联盟或地方响应}

\eventanchor{EH-LEDGER-032-3}{永嘉元年（145年）·冲帝夭亡，梁太后立八岁的刘缵为质帝，梁冀继续掌政}

\eventanchor{EH-LEDGER-032-4}{永嘉元年（145年）·“冲帝夭亡，梁太后立八岁的刘缵为质帝，梁冀继续掌政”的执行中的空间与人员处置}

\eventanchor{EH-LEDGER-032-5}{永嘉元年（145年）·“冲帝夭亡，梁太后立八岁的刘缵为质帝，梁冀继续掌政”的后续制度或军政调整}

\eventanchor{EH-LEDGER-032-6}{永嘉元年（145年）·“冲帝夭亡，梁太后立八岁的刘缵为质帝，梁冀继续掌政”的异源材料与影响范围的复核}

\eventanchor{EH-LEDGER-033-1}{本初元年（146年）·“质帝被毒杀后，梁太后立刘志为桓帝，梁冀继续控制朝政”的前期动员与资源准备}

\eventanchor{EH-LEDGER-033-2}{本初元年（146年）·“质帝被毒杀后，梁太后立刘志为桓帝，梁冀继续控制朝政”的命令、联盟或地方响应}

\eventanchor{EH-LEDGER-033-3}{本初元年（146年）·质帝被毒杀后，梁太后立刘志为桓帝，梁冀继续控制朝政}

\eventanchor{EH-LEDGER-033-4}{本初元年（146年）·“质帝被毒杀后，梁太后立刘志为桓帝，梁冀继续控制朝政”的执行中的空间与人员处置}

\eventanchor{EH-LEDGER-033-5}{本初元年（146年）·“质帝被毒杀后，梁太后立刘志为桓帝，梁冀继续控制朝政”的后续制度或军政调整}

\eventanchor{EH-LEDGER-033-6}{本初元年（146年）·“质帝被毒杀后，梁太后立刘志为桓帝，梁冀继续控制朝政”的异源材料与影响范围的复核}
